\newpage
\phantomsection
\label{prf:fv-msml-distance-to-set-le-distance-to-point-of-mem}
\LRAProofFor{thm:fv-msml-distance-to-set-le-distance-to-point-of-mem}

\begin{remark*}[Return]
Return to \hyperref[thm:fv-msml-distance-to-set-le-distance-to-point-of-mem]{Point-to-Set Distance Is Bounded by Any Witness}.
\end{remark*}

\begin{theorem*}[Point-to-Set Distance Is Bounded by Any Witness]
Let $(X,d)$ be a metric space, let \(x\in X\), and let \(a\in A\). Then
\[
  d(x,A)\leq d(x,a).
\]
\end{theorem*}

\begin{proof}[Professional Standard Proof]
\LRAProofBodyStart
Since \(a\in A\), the number \(d(x,a)\) belongs to the distance set
\[
  \{d(x,y)\mid y\in A\}.
\]
The distance \(d(x,A)\) is the infimum of this set. An infimum is a lower bound
for the set, so it is less than or equal to every member of the set. Applying
this to the member \(d(x,a)\) gives \(d(x,A)\leq d(x,a)\).
\end{proof}

\begin{proof}[Detailed Learning Proof]
\LRAProofBodyStart
The Lean proof starts from the membership hypothesis
\texttt{point\_in\_set : a in A} and extracts two separate facts from it.

First, it constructs a witness proving that \(A\) is nonempty:
\[
  A\neq\emptyset
\]
with witness \(a\) and proof \(a\in A\). In Lean this is the pair
\[
  \langle a,\ \texttt{point\_in\_set}\rangle.
\]
This nonemptiness hypothesis is needed because the lemma
\texttt{distanceToSet\_isGLB} only identifies \(d(x,A)\) as an infimum when
the set \(A\) has at least one point.

Second, it proves that the particular number \(d(x,a)\) belongs to the
distance set from \(x\) to \(A\). Since
\[
  D_{x,A}
  =
  \{d(x,y)\mid y\in A\},
\]
where \(D_{x,A}\) denotes the Lean set \texttt{distanceSet x A},
we again use \(a\) as the witness. The required equality is just
\[
  d(x,a)=d(x,a),
\]
so Lean closes that component with reflexivity. This is the witness package
\[
  \langle a,\ \texttt{point\_in\_set},\ \texttt{rfl}\rangle.
\]

Now invoke \texttt{distanceToSet\_isGLB} with the nonemptiness witness. It
returns
\[
  d(x,A)\ \text{is the greatest lower bound of}\ D_{x,A}.
\]
The first component of an \texttt{IsGLB} proof says that the greatest lower
bound is a lower bound. Applying that lower-bound component to the member
\(d(x,a)\in D_{x,A}\) gives
\[
  d(x,A)\leq d(x,a).
\]
\end{proof}

\begin{proof}[Formalized Lean Proof]
\begin{verbatim}
theorem distanceToSet_le_distance_to_point_of_mem
    {X : Type u}
    [MetricSpace X]
    (x : X)
    {A : Set X}
    {a : X}
    (point_in_set : a in A) :
    distanceToSet x A <= dist x a := by
  have A_nonempty : A.Nonempty := <a, point_in_set>
  have distance_value_in_set :
      dist x a in distanceSet x A := by
    exact <a, point_in_set, rfl>
  have infimum_property :
      IsGLB (distanceSet x A) (distanceToSet x A) := by
    exact distanceToSet_isGLB x A_nonempty
  exact infimum_property.1 distance_value_in_set
\end{verbatim}
\end{proof}

\begin{remark*}[Proof structure]
The membership assumption \(a\in A\) supplies both nonemptiness of \(A\) and a
specific member of \texttt{distanceSet x A}. The proof then uses the
lower-bound half of the \texttt{IsGLB} statement.
\end{remark*}

\NoLocalDependencies

\clearpage
